% Act 1: the oldest form of not doing work twice; the hash that keys the cache starts the content-identity thread.
\begin{frame}[t]{\ev{\tagO\,\tagM}A cache skips work that a previous run did.}
\vspace{.15cm}
\begin{columns}[T,onlytextwidth]
\begin{column}{.56\textwidth}
\begin{tikzpicture}[x=1cm,y=1cm,every node/.style={align=center,font=\fontsize{11.5}{14}\selectfont}, lab/.style={font=\fontsize{10.5}{12}\selectfont\bfseries}]
\node[evidence,text width=3.1cm] (k) at (1.6,0) {\textbf{key} $=$ hash(source,\\flags, toolchain)};
\node[evidence,text width=2.55cm] (hit) at (6.0,0) {Reuse the\\stored output};
\node[candidate,text width=3.1cm] (run) at (1.6,-1.9) {Run the step\\in a sandbox};
\node[state,text width=2.55cm] (st) at (6.0,-1.9) {Store output\\under the key};
\draw[flow] (k)--node[above,lab,text=Teal]{HIT}(hit);
\draw[flow] (k)--node[right,lab,text=Gold]{MISS}(run);
\draw[flow] (run)--(st);
\draw[message] (st)--node[right,lab,text=Teal,align=left]{next\\run}(hit);
\end{tikzpicture}
\par\vspace{.2cm}
{\fontsize{11.5}{14}\selectfont The sandbox hides undeclared files,\\so a step cannot read what the key omits.\par}
\end{column}
\begin{column}{.40\textwidth}
{\fontsize{9.5}{11}\selectfont\color{Teal}\bfseries TOKIO COMPILE\par}
\vspace{.2cm}
\begin{tikzpicture}[x=1cm,y=1cm,every node/.style={font=\fontsize{11}{13}\selectfont}]
\foreach \blabel/\seconds/\y/\col in {Fresh/46/0/Gold,Copied in/13/-.62/Teal,Warm/3/-1.24/Ink}{
\node[anchor=east,align=right] at (1.55,\y+.2) {\blabel};
\fill[\col!65] (1.7,\y) rectangle ({1.7+\seconds*.05},\y+.4);
\node[anchor=west,font=\fontsize{12.5}{15}\selectfont\bfseries] at ({1.8+\seconds*.05},\y+.2) {$\sim$\seconds\ s};}
\end{tikzpicture}
\par\vspace{.2cm}
{\fontsize{10.5}{12.5}\selectfont\color{Muted}Copied in: an earlier runner's build files, not its running state.\par}
\end{column}
\end{columns}
\claim{Distributed builds rely on it: one machine reuses another's output.}
\sources{Tokio: Eliaz, \href{https://github.com/tokio-rs/tokio/issues/8200}{issue \#8200}, 8 June 2026; author-reported external Incredibuild demo, approximate compile times.}
\qadetail{The left diagram is the generic content-addressed action cache used by ccache, Bazel remote caching and Incredibuild: the key is a hash of the declared inputs (sources, flags, toolchain); a hit reuses the stored output; a miss runs the step and stores its output under the key. Sandboxing (hermetic execution) hides undeclared files, so a step cannot silently depend on inputs the key omits; time, randomness, network and nondeterministic tools still need declaring or excluding, so the key covers every input only when the step is hermetic. The hash proves identity of inputs, not correctness of outputs. Tokio: restored cache means Incredibuild compiler artifacts persisted across fresh GitHub Actions runners; hot (already warm) cache denotes the already available cache state in the reported demo. Author-reported external Incredibuild GitHub Actions demo, posted by zozo123 in Tokio issue 8200 on 8 June 2026; approximate compile endpoints 46 s fresh, 13 s restored, 3 s hot. No controlled raw benchmark protocol is posted, and no maintainer adoption is claimed. The issue does not publish enough configuration to isolate disk artifacts, memory/page cache, cache lookup, transfer or incremental checking costs; do not attribute the 13-to-3-second difference entirely to download time or OS memory. Full test path approximately 75 s fresh and 43 s restored; Cargo after touching source approximately 38 s versus 13 s restored Incredibuild. These are static artifact-cache endpoints, not memory-checkpoint or fork timings; the slide labels the issue's restored cache as copied in so that it is not read as restoring execution state, which slide 17 prices. Bazel's local sandbox hides undeclared workspace files, but system directories typically stay visible read-only; hermetic toolchains or a pinned remote-execution image close that gap. We do not claim that a fork is faster than a good build cache.}
\note{[1:05] Start with the oldest way to avoid doing work twice: a cache. Hash a build step's inputs, the source, the flags and the toolchain, into a key. If the key is stored, that is a hit: reuse the stored output. On a miss, run the step in a sandbox and store its output under the key. The sandbox hides files the step did not declare, so it cannot read what the key leaves out. This is the Incredibuild lesson. In my Tokio demo, a fresh compile took about forty-six seconds, copying in an earlier runner's build files about thirteen, and an already warm cache about three. Distributed builds rely on this: one machine reuses another's output. Remember that hash; it comes back.}
\end{frame}
